\section{New Offices and Constitutional Bodies}

\subsection{Constitutional Chamber} \label{subsec:constitutional_chamber}

The Constitutional Chamber is needed to compensate for the deviations from proportionality caused by majority-assured proportional representation.
It is the Bundestag in its composition after the first round of voting and the only institution that can amend the constitution.\footnote{One can discuss whether the Bundesrat should still have to approve all constitutional amendments. For now, I assume that this remains the case.}
If the election was decided after only one round, the Constitutional Chamber is identical to the Bundestag.
Otherwise, it convenes in the composition the Bundestag would have if there had been only one round.
It then also includes candidates who would have been elected after the first round but do not sit in the Bundestag.
It only convenes for its decisions.

For details, see the \mvref{subsec:constitutional-amendments}.

\subsection{Federal Presidential Council} \label{subsec:presidential_council}

It is in the nature of things that constitutional bodies elected by the citizens are particularly susceptible to populism.
Constitutional bodies that are not elected by the citizens can also be infiltrated over time by gradually filling them with enemies of democracy.
But such processes take time.
That is why non-elected constitutional bodies can set limits to populism for a while.

The idea is to let the citizens elect the shaping constitutional bodies (executive and legislature) themselves,
but to have their work controlled by constitutional bodies that are not elected by the citizens.
The controlling constitutional bodies may not themselves shape how the country develops.
After all, they have not been elected.
But they are to check that the elected shapers do their work properly.

We already have one controlling constitutional body.
The Federal Constitutional Court conducts judicial review of legislation.
However, it lacks means of enforcement, and it does not take decisions on matters of state, but only determines what the law is.

I propose strengthening the office of the Federal President in terms of personnel and giving it more authority.
The Federal President should preside over a \enquote{Federal Presidential Council}\footnote{\enquote{State Council} would also be a fitting term. The Federal President as head of state presides over the State Council.
However, the term already carries historical baggage.},
consisting of the President and six other \enquote{Presidential Councillors}.

The Federal President chairs the meetings and decides what goes on the agenda, but otherwise takes decisions on an equal footing with the other members of the Council.

The additional personnel is needed to live up to the increased responsibility for matters of state that is intended for this office.
The Federal Presidential Council is to be able to dissolve the Bundestag or remove the Chancellor in more situations than the Federal President alone can today.
In addition, it should receive veto rights over executive and legislative decisions.

The Federal President retains the representative duties.
The function of head of state, however, is to be shared among seven councillors, including the President.

The Federal Presidential Council is, so to speak, a \enquote{meta-government}.
It does not govern itself, but it installs the government on behalf of the people and can also remove it again if necessary.

\subsubsection{Composition of the Federal Presidential Council} \label{subsubsec:presidential_council_composition}

A controlling constitutional body must be independent of those it is supposed to control.
How to staff institutions that are obliged to be neutral, such as constitutional and supreme courts, offices of the prosecutor general or central banks, with independent people is a long-known problem of democracy\footnote{E.g.\ Bruce Ackerman, \enquote{The New Separation of Powers}, Harvard Law Review 113 (2000), p.~633 ff.; Tarunabh Khaitan, \enquote{Guarantor Institutions}, Asian Journal of Comparative Law 16 (2021), p.~S40 ff.; Mark Tushnet, \enquote{The New Fourth Branch: Institutions for Protecting Constitutional Democracy}, Cambridge University Press 2021.}.
In recent years, there has been growing discussion of how to protect such institutions from being captured or blocked by populists and extremists\footnote{On capture, e.g.\ David Landau, \enquote{Abusive Constitutionalism}, UC Davis Law Review 47 (2013), p.~189 ff.; Kim Lane Scheppele, \enquote{Autocratic Legalism}, University of Chicago Law Review 85 (2018), p.~545 ff.; Steven Levitsky, Daniel Ziblatt, \enquote{How Democracies Die}, Crown 2018.
On blockades, e.g.\ Venice Commission, CDL-PI(2025)023, \url{https://www.coe.int/en/web/venice-commission/-/cdl-pi-2025-023-e}; Gertrude Lübbe-Wolff, \enquote{How to Prevent Blockage of Judicial Appointments}, Verfassungsblog, October~7, 2022, \url{https://verfassungsblog.de/how-to-prevent-blockage-of-judicial-appointments/}; Carolin Lerch, Friedrich Zillessen, \enquote{Parlamentarische Obstruktion}, Verfassungsblog, January~29, 2024, \url{https://verfassungsblog.de/parlamentarische-obstruktion/}.}.
The results of this discussion have already found their way into the 2024 reform of the election of judges to the Federal Constitutional Court.
It can serve as a model, even though it still needs improvement.
How the Federal Presidential Council and the Federal Constitutional Court can be better protected against capture and blockade than before deserves a discussion of its own, which, however, would go beyond the scope of this paper.

\subsection{Federal Office for Opinion Research}\label{subsec:opinion_research_office}

For every government, it is important to know what the people think.
Democratic governments are no exception.
For the proposal made in \cref{subsec:emergency_brake} to take the will of the people into account when deciding on early elections,
a reliable picture of public sentiment is essential.
For such decisions, one cannot rely solely on surveys by private polling institutes.

I therefore propose establishing a Federal Office for Opinion Research that is constantly engaged, without using intelligence-service methods,
in compiling an anonymized picture of public opinion on specific questions and keeping it up to date.
In doing so, it should rely on a larger data basis than the inexpensive representative surveys of private polling institutes.
Their results, however, should definitely also be used as a cross-check.

The Federal Office for Opinion Research is professionally independent, similar to the Federal Statistical Office.
The Office discloses its methodology.

Surveys can be commissioned by the Federal Government, the parliamentary groups of the Bundestag and the Federal Presidential Council.
Government and opposition must have equal access to the Office, not just the government alone.

The commissioning party decides whether and when the results of its survey are published.
There are two reasons for this:
\begin{enumerate}[label=\alph*)]
    \item Government and opposition often want to know what the people think about a question before they take a position or communicate their decision.
    \item Published surveys can trigger an amplification effect, along the lines of: \enquote{If X\% of the population see it that way, then it can't be wrong.}
    If this concerns populist ideas that one wants to counter, one makes life harder for oneself.
    Figures from a state office carry more weight than those from a private institute.
\end{enumerate}

So that the commissioning party cannot cherry-pick the convenient figures, the Office publishes the results itself, and in full, with the wording of the questions, all answers and the methodology.
